% R76 author development: this file extends, and does not replace, the R75 proof.
\section{The complete mean--contrast law and optimal-factor rigidity}
\label{sec:r76-mean-contrast}

The mean of a response and the size of its nonconstant part impose distinct
positivity constraints. The return family of
Theorem~\ref{thm:r75-symmetric-return} fixes the mean at $1/n$.
We now allow every feasible mean and both signs of contrast, for every
$n\ge5$. A second positivity inequality changes the sharp boundary: simply
rescaling the old threshold gives the wrong answer on an interval of means.
The proof yields the complete boundary factorization on that interval.
Throughout, the mean is the average terminal probability over the group;
it is not a prior on an unknown parameter.

Bounded simplex-structured factorization and stacking complementary
nonnegative slacks with a common factor are established techniques
\cite[Section IV-C]{VuThanhGillisLecron23}. Threshold attainment and
boundary identifiability also occur in the geometric theory of nonnegative
rank \cite[Section 4.3]{Gillis12}. Our additional conclusion is the sharp
simultaneous coefficient extremum below, its three-piece parameter region,
and the resulting unrestricted group-memory and boundary-factor laws.

The following elementary extremum treats two nonnegative coefficient
matrices sharing a left null vector. Both outcomes must be represented
through the same preparation; this common-null-vector requirement is
stronger than separate rank bounds on the two outcome matrices.

\begin{lemma}[Sharp simultaneous positivity inequality]
\label{lem:r76-simultaneous-positivity}
Let $n\ge5$ and $v,c,d\in\mathbb R^n$ satisfy
\[
 \sum_i v_i=1,\qquad v\cdot c=v\cdot d=0,\qquad
 c_i\ge v_j-\delta_{ij},\quad d_i\ge\delta_{ij}-v_j
 \quad(1\le i,j\le n).
\]
Put $C=\sum_i c_i$, $D=\sum_i d_i$, and
$\gamma_n=2(n-2)/(n-4)$. For every real $\gamma\ge\gamma_n$,
\begin{equation}\label{eq:r76-weighted-positive-bound}
 \gamma C+D\ \ge\ L_n(\gamma):=
 \frac{\gamma(n+1)+n^2-2n-1+
             2\sqrt{n(\gamma-1)(\gamma+n-2)}}{n-1}.
\end{equation}
Equality holds exactly for the following triples, up to a simultaneous
coordinate permutation. Set
\[
 m=\sqrt{\frac{\gamma+n-2}{n(\gamma-1)}},\qquad
 M=\frac{1+m}{n-1}.
\]
Then
\begin{equation}\label{eq:r76-weighted-equality}
 \begin{split}
 v&=(M,\ldots,M,-m),\\
 c&=(M,\ldots,M,M+M/m),\\
 d&=(1-M,\ldots,1-M,(n-1)M(1-M)/m).
 \end{split}
\end{equation}
In particular $1/\sqrt n<m\le(n-2)/n$. No sign assumption on $v$
or symmetry assumption on the resulting coefficient matrices is made.
\end{lemma}
\begin{proof}
The diagonal inequalities give $D\ge n-1$. If $v\ge0$, the argument
in Lemma~\ref{lem:r75-singular-coefficient} applies verbatim to $c$:
a largest coordinate $M>0$ forces all other positive coordinates to
vanish, hence $C\ge n-1$. We verify below that
$L_n(\gamma)<\gamma(n-1)-1/(n-1)$ in the stated parameter range;
therefore this case is strict.

Otherwise write $m=-\min_i v_i>0$, relabel so that $v_1=M$ is
maximal, let $L$ be the second largest coordinate, and put
$b_0=\max\{L,M-1\}$. The inequalities for $c$ give
\[
 c=(b_0,M,\ldots,M)+u,\quad u\ge0,\qquad
 v\cdot(b_0,M,\ldots,M)=M(b_0+1-M)\ge0.
\]
Using $v\cdot u\ge-m\sum_i u_i$ and $v\cdot c=0$ gives
\begin{equation}\label{eq:r76-C-bound}
 C\ge b_0+(n-1)M+\frac{M(b_0+1-M)}m.
\end{equation}
This is the first-outcome estimate from the retained sharp singular-matrix
proof. For the second outcome, $d=\mathbf1-v+z$ with $z\ge0$ and
$v\cdot z=\|v\|_2^2-1$. Consequently
\[
 D\ge n-1+\frac{1-\|v\|_2^2}{m}.
\]
The bound $-m\le v_i\le M$ implies
$v_i^2\le(M-m)v_i+Mm$; summing it yields
\begin{equation}\label{eq:r76-D-bound}
 D\ge n+\frac{1-M}{m}-nM.
\end{equation}
This bound is valid even when its right side is below $n-1$.

As in the retained proof,
\[
 M\ge M_-:=\frac{1+m}{n-1},\qquad
 L\ge\frac{1-M+m}{n-2},\qquad
 M_+:=1+\frac m{n-1}.
\]
If $M\ge M_+$, \eqref{eq:r76-C-bound} with $b_0\ge M-1$
gives $C\ge nM-1\ge n-1$. Together with $D\ge n-1$ this
case is again strict. If $M_-\le M\le M_+$, substitute
$b_0\ge(1-M+m)/(n-2)$ in \eqref{eq:r76-C-bound}; its right
side increases with $b_0$. This gives
\[
 C\ge F(M):=\frac{1+m}{n-2}+(n-1)M+
                  \frac{n-1}{m(n-2)}(M-M^2).
\]
Therefore
\[
 \gamma C+D\ge G(M):=\gamma F(M)+n+\frac{1-M}{m}-nM.
\]
The quadratic $G$ is strictly concave in $M$, so its minimum on this
interval is at an endpoint. At the left endpoint its value is
\begin{equation}\label{eq:r76-left-endpoint}
 G(M_-)=\frac{\gamma(n+1)+n^2-2n-1+
                  n(\gamma-1)m+(\gamma+n-2)/m}{n-1}
          \ge L_n(\gamma),
\end{equation}
by the arithmetic--geometric mean inequality. At the right endpoint,
\[
 G(M_+)=(\gamma-1)nM_++n-\gamma-\frac1{n-1}
       \ge E_n(\gamma):=\gamma(n-1)-\frac1{n-1}.
\]
For completeness, $E_n(\gamma)>L_n(\gamma)$ throughout the asserted
range. After subtracting the rational terms and multiplying by $n-1$,
this is equivalent to
\[
 n\{(n-3)\gamma-n+2\}
     >2\sqrt{n(\gamma-1)(\gamma+n-2)}.
\]
The left side is positive. On squaring and dividing by $n$, the
difference is
\[
 Q_n(\gamma)=n\{(n-3)\gamma-n+2\}^2
                         -4(\gamma-1)(\gamma+n-2).
\]
For $n\ge6$ its quadratic coefficient is positive,
$Q_n(2)=n(n-6)(n-2)\ge0$, and
$Q_n'(2)=2n(n-3)(n-4)-4n-4>0$. Hence it is strictly positive
at $\gamma\ge\gamma_n>2$. For $n=5$, $\gamma_n=6$ and
$Q_5(\gamma)=16\gamma^2-68\gamma+57$ is positive and increasing
for $\gamma\ge6$. This proves the strict endpoint comparison and
\eqref{eq:r76-weighted-positive-bound}.

Equality forces $M=M_-$ and equality in
\eqref{eq:r76-left-endpoint}; thus $m$ has the stated value. The sum of
the coordinates other than one minimum is $(n-1)M$, so each of them
is $M$. Equality in \eqref{eq:r76-C-bound} places every excess
coordinate of $c$ at the unique negative coordinate and gives
$c_n=M+M/m$. Equality in the bound for $D$ similarly places every
excess over $\mathbf1-v$ at that coordinate, giving the stated $d_n$.
The range $\gamma\ge\gamma_n$ is precisely
$1/\sqrt n<m\le(n-2)/n$, so $M+m\le1$. Thus $d_i=1-M\ge m$
for $i<n$, and $d_n\ge1+m$; all the required inequalities hold.
The displayed $c$ also satisfies its inequalities and $v\cdot c=0$.
Direct substitution verifies $v\cdot d=0$ and equality in the bound.
\end{proof}

Let $\chi_n(g)$ count fixed points in the natural action of $S_n$.
For a mean $p\in[0,1]$ and real contrast $a$, define
\begin{equation}\label{eq:r76-mean-contrast-family}
 \beta_{p,a}(g)=p+a\{\chi_n(g)-1\},\qquad
 b_{p,a}=(\beta_{p,a},1-\beta_{p,a}).
\end{equation}
This is a probability table if and only if
\begin{equation}\label{eq:r76-feasible-mean-contrast}
 0\le p-a\le1,\qquad 0\le p+(n-1)a\le1.
\end{equation}
Both endpoint values occur, at a derangement and at the identity.
Put
\begin{equation}\label{eq:r76-mean-threshold}
 p_0=\frac1{n-\sqrt n},\qquad p_1=\frac2n,\qquad
 F_n(p)=\begin{cases}
 p\,\dfrac{\sqrt n-1}{\sqrt n+1},&0\le p\le p_0,\\[4pt]
 \dfrac{(n-1)p-1}{(n-1)(np-1)},&p_0\le p\le p_1,\\[4pt]
 \dfrac{1-p}{n-1},&p_1\le p\le1.
 \end{cases}
\end{equation}
The three formulas agree at the two meeting points.

\begin{theorem}[Sharp synthesis for bounded centered coefficients]
\label{thm:r76-bounded-coefficients}
Let $n\ge5$, $P=I-J/n$, $0<p<1$ and $a>0$. The following are equivalent:
\begin{enumerate}
\item There are a probability column $u$ and a real matrix $A$ such that
\[
 0\le A\le u\mathbf1^{\mathsf T},\qquad
 \sum_{i,j}A_{ij}=np,\qquad PAP=aP,\qquad
 \operatorname{rank}[A\ \ u]<n.
\]
\item There exists such a pair admitting a common nonnegative factorization
\[
 A=\sum_{i=1}^{n-1}w_i\lambda_i^{\mathsf T}D_i,
 \qquad u=\sum_{i=1}^{n-1}w_i\lambda_i^{\mathsf T},
 \qquad w\in\Delta_{n-1},\quad
 \lambda_i\in\Delta_n,\quad D_i\in[0,1]^n.
\]
\item $a\le F_n(p)$.
\end{enumerate}
Thus joined ordinary rank deficiency and a positive factor of width
$n-1$ have the same parameter-feasibility region in this centered
coefficient problem. This does not assert that every feasible pair
away from the boundary has that nonnegative factorization.

For $0<p\le p_0$ and $a=F_n(p)$, all width-$n-1$ factorizations in (ii), up to factor and simultaneous
coordinate relabeling, are
\begin{equation}\label{eq:r76-low-mean-rigidity}
 \lambda_i=\ell_i,\qquad
 D_i=\frac{np}{(n-1)w_i}\ell_i,\qquad
 \sum_{i=1}^{n-1}w_i=1,\qquad
 w_i\ge\frac{np(1-1/\sqrt n)}{n-1},
\end{equation}
where $\ell_i=(1-1/\sqrt n)e_i+e_n/\sqrt n$ and $D_i$ is the
first-output decoder. For $p_0\le p\le p_1$ and $a=F_n(p)$,
all width-$n-1$ factorizations in (ii) have exactly the form
\begin{equation}\label{eq:r76-middle-rigidity}
 w_i=\frac1{n-1},\qquad
 \lambda_i=(1-b)e_i+b e_n,\qquad
 D_i=e_i+(np-1)e_n,\qquad b=\frac p{np-1},
 \quad 1\le i<n.
\end{equation}
At $p=p_0$ these classifications agree. In the middle interval the
coefficient pair $(A,u)$ itself is unique up to simultaneous coordinate
permutation; the proof gives it explicitly.
\end{theorem}
\begin{proof}
Start with any matrix pair in (i); no positive factorization is assumed
in the converse. Put $B=u\mathbf1^{\mathsf T}-A\ge0$.
A nonzero common left null vector $v$ of $A$ and $u$ also
annihilates $B=u\mathbf1^{\mathsf T}-A$. If its coordinate sum were
zero, $v^{\mathsf T}AP=a v^{\mathsf T}$ would contradict
$v^{\mathsf T}A=0$. Normalize that sum to one.
Exactly as in the retained null-vector decomposition,
\[
 A=a\{I+c\mathbf1^{\mathsf T}-\mathbf1v^{\mathsf T}\},\quad
 B=a\{-I+d\mathbf1^{\mathsf T}+\mathbf1v^{\mathsf T}\},
 \quad v\cdot c=v\cdot d=0.
\]
Nonnegativity gives all the inequalities of
Lemma~\ref{lem:r76-simultaneous-positivity}, and
\begin{equation}\label{eq:r76-mean-null-sums}
 C=\frac pa,\qquad D=\frac{1-p}{a}.
\end{equation}
For $p\le p_0$, apply Lemma~\ref{lem:r75-singular-coefficient} to
$A/(np)$: its centered coefficient is $a/(np)$, so its parameter
there is $a/p$, giving $a\le p\tau_n=F_n(p)$.
For $p\ge p_1$, the diagonal inequalities for $d$ give $D\ge n-1$,
hence $a\le(1-p)/(n-1)=F_n(p)$.

For $p_0<p\le p_1$, set
\[
 m_0=(n-1)p-1,\qquad
 \gamma=\frac{n m_0^2+n-2}{n m_0^2-1}.
\]
Then $1/\sqrt n<m_0\le(n-2)/n$, so
$\gamma\ge\gamma_n$, and $m_0$ is the minimizing value in
\eqref{eq:r76-left-endpoint}. The equality triple
\eqref{eq:r76-weighted-equality} with $M=p$ has sums
\[
 C_0=p(n+1/m_0),\qquad D_0=(1-p)(n+1/m_0).
\]
Thus $L_n(\gamma)=(\gamma p+1-p)(n+1/m_0)$.
Substituting \eqref{eq:r76-mean-null-sums} into the lemma and cancelling
the positive factor $\gamma p+1-p$ gives
\[
 a\le\frac1{n+1/m_0}
     =\frac{(n-1)p-1}{(n-1)(np-1)}=F_n(p).
\]
This proves the necessary boundary everywhere.

We give matching constructions at that boundary. In each case, smaller
$a\ge0$ is obtained on the same core by mixing the displayed decoder
with the constant decoder $p$, with mixing weight $a/F_n(p)$.
For $0<p\le p_0$, use $n-1$ factors with equal weights,
preparation rows $\ell_i$ and decoders $np\ell_i$. The largest decoder
entry is $np(1-1/\sqrt n)\le1$, precisely because $p\le p_0$.
The coefficient matrix is $np A_*$, where $A_*$ is the retained
sharp singular extremizer. Its mean is $p$ and its centered part is
$p\tau_nP=F_n(p)P$, as required.

For $p_0\le p\le p_1$, use \eqref{eq:r76-middle-rigidity}.
Here $0<b<1$ and $0<np-1\le1$. Put $r=np-1$ and
$a_*=(1-b)/(n-1)=F_n(p)$. The resulting coefficient matrix has
\[
 A_{ii}=a_*\ (i<n),\quad A_{ij}=0\ (i\ne j<n),\quad
 A_{in}=a_*r,\quad A_{ni}=\frac b{n-1},\quad A_{nn}=br=p.
\]
The identity $r(nb-1)=1$ shows that all cross differences of
$A-a_*I$ vanish. Thus $PA P=a_*P$ and the group mean is $p$;
the response is $\beta_{p,a_*}$.

For $p_1\le p<1$, put $b=2/n$, use the same preparations
$(1-b)e_i+b e_n$ and equal orbit weights, and let
\[
 d_0=\frac{np-2}{n-2},\qquad
 D_i(j)=d_0+(1-d_0)\mathbf1_{\{j\in\{i,n\}\}}.
\]
These are valid decoders. At $p=p_1$ they agree with the preceding
construction; mixing those decoders with the constant one shows that
the response is $p+[(1-p)/(n-1)](\chi_n-1)$ for every $p\ge p_1$.
Equivalently this follows by direct permutation summation. This proves
attainment at $F_n(p)$. The same
constructions and decoder mixing give (iii)$\Rightarrow$(ii), while
the displayed factorization immediately gives (ii)$\Rightarrow$(i).
The converse already proved (i)$\Rightarrow$(iii), so the matrix
synthesis equivalence is complete.

For the first rigidity statement, the equality case of the retained
singular-coefficient lemma forces $A=np A_*$. The positive diagonal
leading block and its zero off-diagonal entries force precisely one
nonnegative rank-one summand at each of its $n-1$ diagonal positions.
The adjacent last-row and last-column ratios then give
\eqref{eq:r76-low-mean-rigidity}, exactly as in the retained proof.
Its decoder bound is $w_i\ge np(1-1/\sqrt n)/(n-1)$.
At $p=p_0$ summing these lower bounds forces equal weights. Pair
uniqueness at this endpoint also follows without presupposing a positive
factorization. The singular lemma fixes $A=np_0 A_*$, whose normalized
left null vector is $(p_0,\ldots,p_0,-1/\sqrt n)$. Thus
$v\cdot u=0$ and $\sum_i u_i=1$ force $u_n=1/\sqrt n$.
The bounds $u_i\ge A_{ii}=1/(n+\sqrt n)$ for $i<n$ then sum to
$1-1/\sqrt n$, forcing equality in every one of them.

If $p_0<p\le p_1$ and $a=F_n(p)$, equality holds in the weighted
lemma with the chosen $\gamma$. It forces
$v=(p,\ldots,p,-m_0)$ and the unique $c,d$ in
\eqref{eq:r76-weighted-equality}. The coefficient matrix is therefore
exactly the one displayed in the middle construction, while
$u=a(c+d)$ has $u_i=a$ for $i<n$ and $u_n=b$. Its leading
block again forces one rank-one summand at each diagonal position, with
both factors supported on $\{i,n\}$. The last-row ratios force
$\lambda_i=(1-b)e_i+b e_n$. The diagonal entries force
$D_i(i)=a/\{w_i(1-b)\}=1/\{(n-1)w_i\}$.
The upper bound $D_i(i)\le1$ implies $w_i\ge1/(n-1)$;
summing the weights makes every one equal to $1/(n-1)$.
The last-column ratios then give $D_i(n)=np-1$.
This proves the complete middle-boundary classification.
\end{proof}

\begin{corollary}[Exact memory throughout the mean--contrast plane]
\label{thm:r76-mean-contrast}
Let $n\ge5$ and let $(p,a)$ satisfy
\eqref{eq:r76-feasible-mean-contrast}. At $a=0$, both mixed-core and
pure-start minima are one. If $a\ne0$, put $q=p$ when $a>0$ and
$q=1-p$ when $a<0$. Then
\begin{equation}\label{eq:r76-full-memory-law}
 \kappa_{S_n}(b_{p,a})=
 \begin{cases}
 n(n-1),&0<|a|\le F_n(q),\\
 n^2,&|a|>F_n(q),
 \end{cases}
 \qquad m(b_{p,a})=1+\kappa_{S_n}(b_{p,a}).
\end{equation}
Every nonconstant table has observable group $S_n$, residual cardinality
$n!$, and observable dimension $1+(n-1)^2$.
For $a>0$ and $p\ge2/n$, the smaller memory realizes every feasible
contrast. Thus positivity causes no additional core cost anywhere in
that part of the parameter plane.

At the nontrivial boundary $a=F_n(p)$ with $0<p\le2/n$, all
minimum permutation-core factorizations are exactly those in
\eqref{eq:r76-low-mean-rigidity} and \eqref{eq:r76-middle-rigidity}.
For negative contrast, exchange the two output labels. This classifies
minimum permutation cores, without asserting uniqueness of every
unreduced stochastic machine.
\end{corollary}
\begin{proof}
It suffices to treat $a>0$, because exchanging the two outputs replaces
$(p,a)$ by $(1-p,-a)$. The trivial block of the Fourier transform is
$(p,1-p)$, and the nontrivial standard block is
$aI_{n-1}/(n-1)$ in the first coordinate. The faithful-standard-module
argument in Theorem~\ref{thm:r75-symmetric-return} therefore gives the
claimed observable group, residual count and dimension for every $a>0$.
That theorem's subgroup-index estimate depends only on this module.
It implies the lower bound $n(n-1)$ and proves that every core of size
less than $n^2$ compresses to exactly $n-1$ natural orbits: extra orbits
have no standard component, so their nonnegative constant contribution
is absorbed into the natural-orbit decoders. This absorption preserves
both output probabilities and does not use the value $p=1/n$.

On $n$ natural orbits, put weight $1/n$ on each, prepare point $i$
on orbit $i$, and use
\[
 D_i(j)=p-a+na\,\mathbf1_{\{i=j\}}.
\]
These decoders lie in $[0,1]$ exactly under
\eqref{eq:r76-feasible-mean-contrast}. Their response is
\eqref{eq:r76-mean-contrast-family}. Thus the core is always at most
$n^2$, and there are no possible intermediate values.

Suppose a realization on $n-1$ natural orbits exists. Its preparation
rows $\lambda_i$, weights $w_i$ and decoder rows $D_i$ give
\[
 A=\sum_{i<n}w_i\lambda_i^{\mathsf T}D_i,\qquad
 u=\sum_{i<n}w_i\lambda_i^{\mathsf T},\qquad
 B=u\mathbf1^{\mathsf T}-A\ge0.
\]
Here $A\ge0$, $u$ is a probability column, and both coefficient
matrices have the same preparation factor. Their permutation sums give
the two outcomes. The permutation-switch argument in the retained proof
implies $PAP=aP$, and averaging gives $\sum_{ij}A_{ij}=np$.
The preparation matrix has at most $n-1$ rows, so
$\operatorname{rank}[A\ \ u]<n$.
The matrix criterion in
Theorem~\ref{thm:r76-bounded-coefficients} is therefore necessary.
Conversely its width-$n-1$ positive factorizations give exactly
$n-1$ natural-orbit preparations and valid binary decoders.
Their coefficient identities imply the required response for every
permutation. Hence a core below $n^2$ exists exactly when $a\le F_n(p)$,
and then the module lower bound makes its size $n(n-1)$.
The retained pure-start formula gives $m=\min\{n!,1+\kappa\}$,
which equals $1+\kappa$ because $n^2+1<n!$ for $n\ge5$.
At $a=0$ the constant decoder uses one state. Every minimum core on the
nontrivial boundary consists of $n-1$ natural orbits, so the factor
classification follows from the matrix theorem.
\end{proof}

\paragraph{Physical realization and scope.}
Use the prepared unitary experiment of
Theorem~\ref{thm:r75-symmetric-return}, with
$\rho=n^{-1}\sum_i|ii\rangle\langle ii|$ and
$E=\sum_i|ii\rangle\langle ii|$. Replacing the terminal effect by
\[
 E_{p,a}=(p-a)I+naE
\]
produces \eqref{eq:r76-mean-contrast-family}; its two eigenvalues are
$p-a$ and $p+(n-1)a$, so the exact validity range is
\eqref{eq:r76-feasible-mean-contrast}, including negative contrast.
The same two command generators realize every word.
The theorem solves the whole displayed two-parameter family, including
its constant boundary and its minimum-core factorization on the
nontrivial transition curve. It is not an unrestricted high-contrast
formula for every finite group, an identification result for an unknown
model, or a solution of the original arbitrary-prior decision-risk problem.
